\documentclass[11pt]{article}
\input{../../shared/project_style.tex}
\rhead{Basics 22 Textbook Draft}

\begin{document}
\DocTitle{Basics 22: Moments of the Kinetic Equation}{III - Elementary plasma physics}{How continuity, momentum, pressure, and energy equations emerge from phase-space dynamics}

\begin{overviewbox}
\textbf{Textbook draft, version 1.0.} Fluid plasma equations are velocity moments of a kinetic equation. This chapter performs the first three moment derivations carefully, showing where the pressure tensor, heat flux, collision terms, and closure problem enter. This chapter is designed for a reader with no prior plasma-physics training. It distinguishes exact identities from physical approximations and connects the central equations to later simulation modules.
\end{overviewbox}

\ProjectTOC

\section{Learning objectives}
After completing this note, the reader should be able to:
\begin{itemize}[leftmargin=1.6em]
\item Define velocity moments of a distribution function.
\item Derive continuity from the kinetic equation.
\item Derive momentum balance and the pressure tensor.
\item Derive internal- and total-energy moment equations.
\item Explain the infinite moment hierarchy.
\item Distinguish conservative and convective fluid forms.
\end{itemize}

\section{Prerequisites and reading strategy}
\begin{itemize}[leftmargin=1.6em]
\item Basics 3, 5, 13, and 19.
\end{itemize}
On a first reading, emphasize physical meaning and dimensional checks. On a second reading, reproduce the derivations. Attempt the computational laboratory only after the limiting cases are understood.

\section{Starting kinetic equation}

For species $s$,
\begin{equation}
\frac{\partial f_s}{\partial t}
+\mathbf v\cdot\nabla_{\mathbf x}f_s
+\frac{q_s}{m_s}(\mathbf E+\mathbf v\times\mathbf B)
\cdot\nabla_{\mathbf v}f_s
=C_s[f]+S_s.
\end{equation}
Assume $f_s$ decays sufficiently rapidly as $|\mathbf v|\to\infty$ that velocity-space surface terms vanish. This technical condition is what permits integration by parts.


\section{Zeroth moment: continuity}

Integrate over velocity:
\begin{equation}
\frac{\partial}{\partial t}\int f_s\,d^3v
+\nabla\cdot\int\mathbf v f_s\,d^3v
+\frac{q_s}{m_s}\int(\mathbf E+\mathbf v\times\mathbf B)\cdot\nabla_{\mathbf v}f_s\,d^3v
=\int(C_s+S_s)\,d^3v.
\end{equation}
The force term vanishes by velocity-space integration by parts because $\nabla_{\mathbf v}\cdot(\mathbf E+\mathbf v\times\mathbf B)=0$. Thus
\begin{equation}
\boxed{\frac{\partial n_s}{\partial t}+\nabla\cdot(n_s\mathbf u_s)=S_{n,s}}.
\end{equation}
Number-conserving collisions satisfy $\int C_s\,d^3v=0$.


\section{First moment: momentum}

Multiply by $m_s\mathbf v$ and integrate. The force term becomes
\begin{equation}
q_sn_s(\mathbf E+\mathbf u_s\times\mathbf B).
\end{equation}
Define the momentum-flux tensor
\begin{equation}
m_s\int\mathbf v\mathbf v f_s\,d^3v
=m_sn_s\mathbf u_s\mathbf u_s+\mathsf P_s.
\end{equation}
The conservative momentum equation is
\begin{equation}
\frac{\partial(m_sn_s\mathbf u_s)}{\partial t}
+\nabla\cdot(m_sn_s\mathbf u_s\mathbf u_s+\mathsf P_s)
=q_sn_s(\mathbf E+\mathbf u_s\times\mathbf B)+\mathbf R_s+\mathbf M_s.
\end{equation}
Using continuity converts it to convective form.


\section{Second moment: kinetic and internal energy}

Multiply by $m_sv^2/2$ and integrate:
\begin{equation}
\frac{\partial \mathcal E_s}{\partial t}
+\nabla\cdot\mathbf F_{\mathcal E,s}
=q_sn_s\mathbf u_s\cdot\mathbf E+Q_s.
\end{equation}
The magnetic force performs no work. Decomposing $\mathbf v=\mathbf u_s+\mathbf c$ gives
\begin{equation}
\mathcal E_s=\frac{1}{2}m_sn_su_s^2+\mathcal U_s,\qquad
\mathcal U_s=\frac{1}{2}\operatorname{tr}\mathsf P_s.
\end{equation}
For an isotropic monatomic distribution, $\mathcal U_s=3p_s/2$. The energy flux contains bulk transport, pressure work, and heat flux
\begin{equation}
\mathbf q_s=\frac{m_s}{2}\int c^2\mathbf c f_s\,d^3v.
\end{equation}


\section{Pressure-tensor evolution and the hierarchy}

Multiplying by $m_s\mathbf c\mathbf c$ produces a pressure-tensor equation involving the third central moment, while the heat-flux equation involves a fourth moment. Thus
\begin{equation}
\text{moment }k\quad\longrightarrow\quad\text{moment }k+1.
\end{equation}
No finite set closes itself without additional physics. Collisions may justify isotropic pressure and Fourier-like heat flow; collisionless magnetized plasmas often require anisotropic or Landau-fluid closures.


\section{Conservative versus primitive formulations}

A conservative equation has the form
\begin{equation}
\frac{\partial U}{\partial t}+\nabla\cdot\mathbf F(U)=S(U).
\end{equation}
It naturally represents integral balances and shocks. Primitive equations evolve variables such as $n$, $\mathbf u$, and $p$ and often expose physical interpretation. Numerical solvers may reconstruct primitive variables but update conserved variables to maintain correct balances.


\section{Connection to kinetic MHD}

Ideal MHD closes the bulk hierarchy with scalar pressure and an adiabatic law. Energetic-particle kinetic MHD instead keeps a distribution or selected kinetic moments for the fast population. Its pressure perturbation then feeds the fluid momentum or energy principle, retaining wave-particle resonances that scalar bulk closure cannot produce.


\section{Computational laboratory}
Represent a distribution on a velocity grid and evaluate density, momentum, pressure tensor, energy, and heat flux by quadrature. Verify the moments of a shifted Maxwellian and bi-Maxwellian. Demonstrate how velocity-domain truncation violates conservation.

\section{Summary}
\begin{itemize}[leftmargin=1.6em]
\item Fluid equations are weighted velocity integrals of the kinetic equation.
\item Continuity, momentum, and energy correspond to the first three physically important moments.
\item Pressure and heat flux encode unresolved velocity-space structure.
\item The moment hierarchy is infinite and requires closure.
\item Kinetic MHD preserves selected kinetic moments or distributions where fluid closure is inadequate.
\end{itemize}

\SymbolsAcronymsSection
\begin{symtable}
$f_s$ & distribution function & Phase-space density of species $s$. \\
$\mathsf P_s$ & pressure tensor & Second central moment. \\
$\mathbf q_s$ & heat flux & Third central moment. \\
$\mathcal E_s$ & total particle energy density & Bulk plus internal kinetic energy. \\
$C_s[f]$ & collision operator & Velocity-space interaction term. \\
$S_{n,s}$ & number source & Zeroth moment of external sources. \\
\end{symtable}

\section{Exercises}
\begin{enumerate}[leftmargin=1.8em]
\item Derive the zeroth moment and show why the Lorentz force contributes no particle source.
\item Derive the electromagnetic force in the first moment by integration by parts.
\item Decompose the momentum-flux tensor into bulk and pressure terms.
\item Show that magnetic forces do not appear in the kinetic-energy source.
\item Explain why a scalar equation of state is a closure rather than a consequence of the kinetic equation.
\end{enumerate}

\section{References}
\begin{thebibliography}{99}
\bibitem{ref1} G. B. Arfken, H. J. Weber, and F. E. Harris, \emph{Mathematical Methods for Physicists}, 7th ed., Academic Press (2013).
\bibitem{ref2} G. Strang, \emph{Linear Algebra and Its Applications}, 4th ed., Brooks/Cole (2006).
\bibitem{ref3} W. A. Strauss, \emph{Partial Differential Equations: An Introduction}, 2nd ed., Wiley (2007).
\bibitem{ref4} F. F. Chen, \emph{Introduction to Plasma Physics and Controlled Fusion}, 3rd ed., Springer (2016).
\bibitem{ref5} R. J. Goldston and P. H. Rutherford, \emph{Introduction to Plasma Physics}, CRC Press (1995).
\bibitem{ref6} P. M. Bellan, \emph{Fundamentals of Plasma Physics}, Cambridge University Press (2006).
\bibitem{ref7} J. P. Freidberg, \emph{Ideal MHD}, Cambridge University Press (2014).
\bibitem{ref8} J. P. Goedbloed, R. Keppens, and S. Poedts, \emph{Advanced Magnetohydrodynamics}, Cambridge University Press (2010).
\bibitem{ref9} H. Goedbloed and S. Poedts, \emph{Principles of Magnetohydrodynamics}, Cambridge University Press (2004).
\end{thebibliography}

\end{document}
